\documentclass[12pt]{article}
\usepackage[utf8]{inputenc}
\usepackage[T1]{fontenc}
\usepackage{hyperref}
\title{Tsogo Language Resources for K-12 Curriculum Development}
\author{Compiled by Claude Sonnet 5 (claude-sonnet-5) for Andr\'as Kornai}
\date{}
\begin{document}
\maketitle
\section{Introduction}
Tsogo (ISO 639-3 tsv; Guthrie B.31) is a Bantu language of Gabon spoken mainly by the Mitsogho, with about 9,000 speakers (2007 estimate). The search found very little beyond encyclopedic background.
\section{Primary Educational Resources}
None found.
\section{Linguistic and Reference Materials}
Overview articles: \url{https://en.wikipedia.org/wiki/Tsogo_language} and \url{https://en.wikipedia.org/wiki/Tsogo_languages}. No published Tsogo dictionary or grammar turned up. A Stellenbosch study of encyclopaedic dictionaries for Gabonese languages uses Yilumbu as its case and is only indirectly relevant: \url{https://scholar.sun.ac.za/handle/10019.1/3990}.
\section{Religious and Community Resources}
None found.
\section{Digital and Audio Resources}
None found.
\section{Curriculum Development Context}
Little exists in what the search returned. Any curriculum would need new documentation work. Searches under the names Mitsogho or Getsogo may find more.
\end{document}